\section{Related Work}
\label{sec:related}

\paragraph{Voting Rights Act Section 2 and the Gingles test.}
Section 2 of the Voting Rights Act of 1965, as amended in 1982,
prohibits electoral practices that result in minority voters having less
opportunity than other members to participate in the political process
\citep{vra1965}.
In \emph{Thornburg v.\ Gingles} (1986), the Supreme Court established
a three-part test for Section 2 vote dilution claims in redistricting
\citep{gingles1986}:

\begin{enumerate}
\item \textbf{Prong 1 (geographic compactness)}: The minority group is
  sufficiently large and geographically compact to constitute a majority
  in a single-member district.
\item \textbf{Prong 2 (political cohesion)}: The minority group is
  politically cohesive.
\item \textbf{Prong 3 (bloc voting)}: The white majority votes sufficiently
  as a bloc to usually defeat the minority's preferred candidate.
\end{enumerate}

Gingles Prong 1 is the geographic predicate: it asks whether the minority
community is compact enough to form a district.
VRASection operationalises Prong 1 directly: the alignment score $A$ measures
geographic concentration, and $A \geq 0.15$ (proposed threshold)
identifies communities concentrated enough for VRASection to detect them.

\paragraph{Shaw v. Reno and racial predominance.}
\emph{Shaw v.\ Reno} (1993) and \emph{Miller v.\ Johnson} (1995)
established that racial considerations cannot be the ``predominant factor''
in district design \citep{shaw1993,miller1995}.
A district is presumptively unconstitutional if it is so ``bizarrely shaped''
that it is unexplainable on grounds other than race.
The predominance test asks whether the legislature subordinated traditional
redistricting principles --- compactness, contiguity, preservation of
political subdivisions --- to racial considerations.

VRASection's $w_\text{vra}$ parameter makes the Shaw analysis quantitative:
at $w_\text{vra} < 0.5$, the compactness component of the selection score
exceeds the alignment component for any ratio, ensuring that geographic
compactness is the predominant factor (Section~\ref{sec:legal-analysis}).
At $w_\text{vra} = 0.40$ (the proposed default), compactness is explicitly
weighted 60\,\% against alignment's 40\,\%.

\paragraph{Allen v. Milligan and the modern Section 2 standard.}
\emph{Allen v.\ Milligan} (2023) reaffirmed the Gingles framework and
held that Alabama's 2021 congressional map, which produced only one
Black-majority district despite Alabama's 27\,\% Black VAP, violated
Section 2 \citep{allen2023}.
The Court's 5--4 opinion confirmed that race-neutral algorithms do not
insulate a map from Section 2 challenge if the resulting map dilutes
minority voting strength.
Alabama is the lead case study for VRASection precisely because its
geography --- concentrated Black Belt population, dispersed northern
industrial and coastal areas --- provides a strong alignment signal that
the algorithm should detect and act on.

\paragraph{Callais v. Landry and the Louisiana remedy.}
\emph{Callais v.\ Landry} (2025) addressed Louisiana's remedial maps
following Section 2 litigation \citep{callais2025}.
The case establishes that the remedy --- two Black-majority congressional
districts --- must be achieved while maintaining other traditional
redistricting criteria.
VRASection is designed to produce plans that satisfy this requirement:
the alignment score concentrates minority population without abandoning
compactness or contiguity.
Validation against Louisiana's Section 2 remedy (2 Black-majority districts
from $k=6$) is a key experiment in the full empirical suite.

\paragraph{Geographic concentration and the Rodden thesis.}
\citet{rodden2019} documents the deep geographic sorting of American voters:
Democrats are concentrated in dense urban cores; Republicans are distributed
across suburban and rural areas.
This sorting interacts with redistricting in two ways.
For partisan purposes, minimum-edge-cut bisection tends to isolate urban
cores, producing the caterpillar pathology identified in B.8 \citep{b8paper}.
For VRA purposes, the Rodden thesis implies that minority communities
(disproportionately Democratic) are already geographically concentrated
in urban cores and specific rural corridors (e.g., the Black Belt, the
Rio Grande Valley, the Mississippi Delta).
VRASection exploits this concentration directly: states where minority
communities have high geographic clustering will exhibit high alignment
scores, and the algorithm's tie-breaking will concentrate them on one
side of the first bisection.

\citet{chen2013} confirm that geographic sorting produces predictable
partisan outcomes under compactness criteria.
VRASection's alignment score provides a parallel signal for minority
communities: the geographic sorting that makes a community legally cognizable
under Gingles Prong 1 is the same sorting that produces a high alignment score.

\paragraph{Compactness in the VRA context.}
Compactness serves a dual function in VRA jurisprudence.
Under Gingles Prong 1, compactness is a predicate for coverage: only
compact minority communities trigger Section 2's affirmative obligation.
Under Shaw, compactness is a limiting principle: racial predominance is
established when compactness is sacrificed in service of racial targets.
VRASection navigates this tension by making compactness the primary criterion
(the $\mathrm{EC}_\text{norm}$ term in the objective) and alignment a bounded
secondary term ($w_\text{vra} \leq 0.5$).
The algorithm satisfies Gingles Prong 1's geographic compact requirement
by rewarding natural geographic concentration --- the same geographic
compactness required by Shaw.

\paragraph{MetisVra (B.3): the explicit constraint approach.}
B.3 \citep{b3paper} introduced MetisVra, which adds minority VAP as an
explicit second METIS constraint (ncon=2).
The approach fails in 3 of 5 tested states due to constraint scale mismatch:
minority VAP fractions (0.0--1.0) and census population counts
(10,000--200,000) differ by 60--800$\times$, and METIS's numerical solver
cannot simultaneously satisfy both constraints with meaningful precision.
Alabama (0\,\% success across 140 configurations) and Louisiana (0\,\%)
are the clearest failures; South Carolina also fails completely.
MetisVra succeeds only in states with large $k$ (Georgia, $k=14$) or
very high minority concentration (Mississippi, $k=4$, 38.8\,\% Black VAP).

VRASection addresses the MetisVra failure at the architectural level:
there is no second METIS constraint, so there is no scale mismatch.
The alignment score is computed post-hoc from the METIS output, not
used as a constraint input.

\paragraph{GeoSection (B.8): the isoperimetric ratio scan.}
VRASection extends GeoSection \citep{b8paper}, which introduced
isoperimetrically normalised ratio scanning for recursive bisection.
GeoSection's key contribution --- the $\mathrm{EC}/\sqrt{\min(i,k-i)}$
normalisation that prevents the caterpillar pathology --- carries over
to VRASection unchanged.
VRASection adds a single modification: the selection score at the first
bisection level incorporates the alignment bonus.
All subsequent recursion levels use standard GeoSection (no alignment),
so that race is observed at most once in the bisection tree.
